% =====================================================================
% Section 4 -- Exploratory Feature Analysis (P2)
% For I1: % =====================================================================
% Section 4 -- Exploratory Feature Analysis (P2)
% For I1: % =====================================================================
% Section 4 -- Exploratory Feature Analysis (P2)
% For I1: % =====================================================================
% Section 4 -- Exploratory Feature Analysis (P2)
% For I1: \input{p2_section4} inside the ACL template.
% Requires: graphicx, booktabs. Figures live in EDA/figures/ (PDF and PNG).
% Every number below is produced by `python EDA/p2_eda.py` from the
% TRAINING portion only (EDA/outputs/*.csv). Citation keys to add to the
% .bib: huang2021interpretability, ayresdecampos2000sisporto.
% Section references (\S7--\S8) follow the course outline; adjust if needed.
% =====================================================================
\section{Exploratory Feature Analysis}
\label{sec:eda}

All analyses in this section use only the training portion (972 Normal and
104 Pathologic segments from 175 examinations); validation and test data were
not inspected. Labels are used here for description only; the one-class
detectors are fitted on Normal training segments. Because the 104 Pathologic
segments come from only 28 examinations, and two examinations contribute 36 of
them, uncertainty is reported with a cluster bootstrap that resamples whole
examinations (2{,}000 resamples).

\paragraph{Distribution contrast.}
Figure~\ref{fig:variability} compares the four FHR-variability features.
ASTV, the percentage of time with abnormal short-term variability, shifts
upward (median 40 vs.\ 65), and MLTV collapses toward zero (51\% of Pathologic
segments have $\mathrm{MLTV}=0$ vs.\ 4\% of Normal ones). MSTV is instead
bimodal in the Pathologic class. Splitting Pathologic segments by a label-free
rule explains this: segments with prolonged decelerations ($\mathrm{DP}>0$;
65 segments) have high MSTV, a wide histogram and a low histogram mean (medians
2.7, 113 and 96\,bpm), whereas those without ($\mathrm{DP}=0$; 39 segments)
are flat (medians: MSTV 0.3, Width 18, ALTV 69\%). The two groups deviate from
Normal (MSTV 1.2, Width 68, Mean 134\,bpm, ALTV 0\%) in \emph{opposite}
directions; only ASTV is elevated in both (ASTV $\ge 60$ in 83\% and 95\% of
the two groups vs.\ 16\% of Normal). Event features are zero-inflated, so we
compare presence rates: a prolonged deceleration occurs in 62.5\% of
Pathologic but 4.7\% of Normal segments, and an acceleration in 19.2\% vs.\
67.2\%.

\begin{figure*}[t]
  \centering
  \includegraphics[width=\textwidth]{fig1_variability_contrast.pdf}
  \caption{FHR-variability features in the training portion. Top:
  histograms normalised within each class, so that the 972:104 imbalance does
  not distort the comparison. Bottom: empirical CDFs; the KS statistic is the
  largest vertical gap. ASTV and ALTV are percentages of time; UCI states no
  unit for MSTV and MLTV. Pathologic MSTV is bimodal and 69\% of Pathologic
  ALTV values are zero.}
  \label{fig:variability}
\end{figure*}

\paragraph{Feature separability.}
Table~\ref{tab:sep} reports the two-sample Kolmogorov--Smirnov statistic, a
rank-based AUC, and $k$NN mutual information (MI) with the label. Wasserstein
distances were computed after scaling by the Normal standard deviation, since
raw values are not comparable across units (ASTV: 1.40 SD; Mean: 1.84 SD).
ASTV is the most separable and the most stable feature (KS 0.71, 95\% CI
0.56--0.84). DP and the histogram mean rank next, but their intervals are wide
and their KS drops to 0.41 and 0.37 when the two largest Pathologic
examinations are removed, so they partly reflect a few examinations. A naive
segment-level bootstrap would have given 0.51--0.68 for DP, which understates
this uncertainty. MSTV shows why direction-specific scores mislead for
anomaly detection: its AUC is 0.59, but the AUC of its absolute deviation from
the Normal mean is 0.80. The baseline LB carries no usable signal on its own
(0.37).

\begin{table}[t]
\centering
\small
\setlength{\tabcolsep}{2.8pt}
\begin{tabular}{lrrlccc}
\toprule
 & \multicolumn{2}{c}{Median} & & & & \\
\cmidrule(lr){2-3}
Feature & N & P & KS [95\% CI] & AUC & $|z|$ & MI \\
\midrule
ASTV (\%)     & 40    & 65    & .71 [.56, .84] & .86 & .83 & .13 \\
DP$^{\dagger}$ & 0 & 1.7   & .60 [.29, .78] & .80 & .80 & .13 \\
Mean (bpm)    & 134   & 103.5 & .58 [.29, .77] & .19 & .80 & .14 \\
MLTV          & 8.1   & 0     & .54 [.44, .68] & .17 & .73 & .09 \\
AC$^{\dagger}$ & 2.5 & 0   & .50 [.34, .67] & .23 & .57 & .05 \\
MSTV          & 1.2   & 2.1   & .36 [.31, .59] & .59 & .80 & .13 \\
ALTV (\%)     & 0     & 0     & .22 [.14, .50] & .52 & .63 & .06 \\
LB (bpm)      & 132   & 133   & .22 [.15, .39] & .49 & .37 & .08 \\
\bottomrule
\end{tabular}
\caption{Separability of selected features (training portion; N = Normal,
P = Pathologic). AUC = $P(x_P > x_N)$ with ties counted one half, so values
far from .5 in either direction indicate separation; $|z|$ = AUC of the
absolute $z$-score with respect to the Normal mean and SD; MI in nats (label
entropy 0.32); $\dagger$ events per 1000\,s. CIs resample examinations. All 21 features are in the
supplementary table.}
\label{tab:sep}
\end{table}

\paragraph{Redundancy.}
The histogram family is strongly redundant. $\mathrm{Width}=\mathrm{Max}-\mathrm{Min}$
holds exactly for every training segment, so the covariance of the 21 features
has rank 20. Mode, Mean and Median are nearly interchangeable (Spearman
$\rho=0.90$--$0.97$), and Width, Min and Variance are strongly related
($\rho=-0.89$ and $0.82$). In contrast, the variability features are not
redundant with one another (ASTV--MSTV $\rho=-0.36$; ALTV--MLTV $\rho=-0.02$).
A label-free rule fitted on Normal training segments (complete-linkage
clusters with $|\rho|\ge 0.8$, keep the medoid, drop Width) removes Width,
Min, Mode and Mean, leaving 17 features.

\paragraph{Exploratory hypotheses.}
These observations shape the feature design tested in \S7--\S8. Pathologic
segments form two phenotypes that are visible in different feature families,
which motivates family-level rather than single-feature ablations.
\textbf{H1}: the variability family (ASTV, MSTV, ALTV, MLTV) is necessary;
removing it lowers detection performance, mainly on Pathologic segments
without prolonged decelerations.
\textbf{H2}: deceleration type is informative; replacing DL, DS and DP by
their sum does not help, because light decelerations are common in Normal
segments (45\%) and the sum is less separable than DP alone (KS 0.40 vs.\
0.60). Our initial expectation that aggregated deceleration features would
improve on the raw ones was therefore not supported.
\textbf{H3}: the 17-feature set obtained by the label-free pruning rule
performs on par with all 21 features. Each hypothesis changes only the feature
set, with the data split, model configuration and seeds held fixed, and is
judged on validation PR-AUC.
 inside the ACL template.
% Requires: graphicx, booktabs. Figures live in EDA/figures/ (PDF and PNG).
% Every number below is produced by `python EDA/p2_eda.py` from the
% TRAINING portion only (EDA/outputs/*.csv). Citation keys to add to the
% .bib: huang2021interpretability, ayresdecampos2000sisporto.
% Section references (\S7--\S8) follow the course outline; adjust if needed.
% =====================================================================
\section{Exploratory Feature Analysis}
\label{sec:eda}

All analyses in this section use only the training portion (972 Normal and
104 Pathologic segments from 175 examinations); validation and test data were
not inspected. Labels are used here for description only; the one-class
detectors are fitted on Normal training segments. Because the 104 Pathologic
segments come from only 28 examinations, and two examinations contribute 36 of
them, uncertainty is reported with a cluster bootstrap that resamples whole
examinations (2{,}000 resamples).

\paragraph{Distribution contrast.}
Figure~\ref{fig:variability} compares the four FHR-variability features.
ASTV, the percentage of time with abnormal short-term variability, shifts
upward (median 40 vs.\ 65), and MLTV collapses toward zero (51\% of Pathologic
segments have $\mathrm{MLTV}=0$ vs.\ 4\% of Normal ones). MSTV is instead
bimodal in the Pathologic class. Splitting Pathologic segments by a label-free
rule explains this: segments with prolonged decelerations ($\mathrm{DP}>0$;
65 segments) have high MSTV, a wide histogram and a low histogram mean (medians
2.7, 113 and 96\,bpm), whereas those without ($\mathrm{DP}=0$; 39 segments)
are flat (medians: MSTV 0.3, Width 18, ALTV 69\%). The two groups deviate from
Normal (MSTV 1.2, Width 68, Mean 134\,bpm, ALTV 0\%) in \emph{opposite}
directions; only ASTV is elevated in both (ASTV $\ge 60$ in 83\% and 95\% of
the two groups vs.\ 16\% of Normal). Event features are zero-inflated, so we
compare presence rates: a prolonged deceleration occurs in 62.5\% of
Pathologic but 4.7\% of Normal segments, and an acceleration in 19.2\% vs.\
67.2\%.

\begin{figure*}[t]
  \centering
  \includegraphics[width=\textwidth]{fig1_variability_contrast.pdf}
  \caption{FHR-variability features in the training portion. Top:
  histograms normalised within each class, so that the 972:104 imbalance does
  not distort the comparison. Bottom: empirical CDFs; the KS statistic is the
  largest vertical gap. ASTV and ALTV are percentages of time; UCI states no
  unit for MSTV and MLTV. Pathologic MSTV is bimodal and 69\% of Pathologic
  ALTV values are zero.}
  \label{fig:variability}
\end{figure*}

\paragraph{Feature separability.}
Table~\ref{tab:sep} reports the two-sample Kolmogorov--Smirnov statistic, a
rank-based AUC, and $k$NN mutual information (MI) with the label. Wasserstein
distances were computed after scaling by the Normal standard deviation, since
raw values are not comparable across units (ASTV: 1.40 SD; Mean: 1.84 SD).
ASTV is the most separable and the most stable feature (KS 0.71, 95\% CI
0.56--0.84). DP and the histogram mean rank next, but their intervals are wide
and their KS drops to 0.41 and 0.37 when the two largest Pathologic
examinations are removed, so they partly reflect a few examinations. A naive
segment-level bootstrap would have given 0.51--0.68 for DP, which understates
this uncertainty. MSTV shows why direction-specific scores mislead for
anomaly detection: its AUC is 0.59, but the AUC of its absolute deviation from
the Normal mean is 0.80. The baseline LB carries no usable signal on its own
(0.37).

\begin{table}[t]
\centering
\small
\setlength{\tabcolsep}{2.8pt}
\begin{tabular}{lrrlccc}
\toprule
 & \multicolumn{2}{c}{Median} & & & & \\
\cmidrule(lr){2-3}
Feature & N & P & KS [95\% CI] & AUC & $|z|$ & MI \\
\midrule
ASTV (\%)     & 40    & 65    & .71 [.56, .84] & .86 & .83 & .13 \\
DP$^{\dagger}$ & 0 & 1.7   & .60 [.29, .78] & .80 & .80 & .13 \\
Mean (bpm)    & 134   & 103.5 & .58 [.29, .77] & .19 & .80 & .14 \\
MLTV          & 8.1   & 0     & .54 [.44, .68] & .17 & .73 & .09 \\
AC$^{\dagger}$ & 2.5 & 0   & .50 [.34, .67] & .23 & .57 & .05 \\
MSTV          & 1.2   & 2.1   & .36 [.31, .59] & .59 & .80 & .13 \\
ALTV (\%)     & 0     & 0     & .22 [.14, .50] & .52 & .63 & .06 \\
LB (bpm)      & 132   & 133   & .22 [.15, .39] & .49 & .37 & .08 \\
\bottomrule
\end{tabular}
\caption{Separability of selected features (training portion; N = Normal,
P = Pathologic). AUC = $P(x_P > x_N)$ with ties counted one half, so values
far from .5 in either direction indicate separation; $|z|$ = AUC of the
absolute $z$-score with respect to the Normal mean and SD; MI in nats (label
entropy 0.32); $\dagger$ events per 1000\,s. CIs resample examinations. All 21 features are in the
supplementary table.}
\label{tab:sep}
\end{table}

\paragraph{Redundancy.}
The histogram family is strongly redundant. $\mathrm{Width}=\mathrm{Max}-\mathrm{Min}$
holds exactly for every training segment, so the covariance of the 21 features
has rank 20. Mode, Mean and Median are nearly interchangeable (Spearman
$\rho=0.90$--$0.97$), and Width, Min and Variance are strongly related
($\rho=-0.89$ and $0.82$). In contrast, the variability features are not
redundant with one another (ASTV--MSTV $\rho=-0.36$; ALTV--MLTV $\rho=-0.02$).
A label-free rule fitted on Normal training segments (complete-linkage
clusters with $|\rho|\ge 0.8$, keep the medoid, drop Width) removes Width,
Min, Mode and Mean, leaving 17 features.

\paragraph{Exploratory hypotheses.}
These observations shape the feature design tested in \S7--\S8. Pathologic
segments form two phenotypes that are visible in different feature families,
which motivates family-level rather than single-feature ablations.
\textbf{H1}: the variability family (ASTV, MSTV, ALTV, MLTV) is necessary;
removing it lowers detection performance, mainly on Pathologic segments
without prolonged decelerations.
\textbf{H2}: deceleration type is informative; replacing DL, DS and DP by
their sum does not help, because light decelerations are common in Normal
segments (45\%) and the sum is less separable than DP alone (KS 0.40 vs.\
0.60). Our initial expectation that aggregated deceleration features would
improve on the raw ones was therefore not supported.
\textbf{H3}: the 17-feature set obtained by the label-free pruning rule
performs on par with all 21 features. Each hypothesis changes only the feature
set, with the data split, model configuration and seeds held fixed, and is
judged on validation PR-AUC.
 inside the ACL template.
% Requires: graphicx, booktabs. Figures live in EDA/figures/ (PDF and PNG).
% Every number below is produced by `python EDA/p2_eda.py` from the
% TRAINING portion only (EDA/outputs/*.csv). Citation keys to add to the
% .bib: huang2021interpretability, ayresdecampos2000sisporto.
% Section references (\S7--\S8) follow the course outline; adjust if needed.
% =====================================================================
\section{Exploratory Feature Analysis}
\label{sec:eda}

All analyses in this section use only the training portion (972 Normal and
104 Pathologic segments from 175 examinations); validation and test data were
not inspected. Labels are used here for description only; the one-class
detectors are fitted on Normal training segments. Because the 104 Pathologic
segments come from only 28 examinations, and two examinations contribute 36 of
them, uncertainty is reported with a cluster bootstrap that resamples whole
examinations (2{,}000 resamples).

\paragraph{Distribution contrast.}
Figure~\ref{fig:variability} compares the four FHR-variability features.
ASTV, the percentage of time with abnormal short-term variability, shifts
upward (median 40 vs.\ 65), and MLTV collapses toward zero (51\% of Pathologic
segments have $\mathrm{MLTV}=0$ vs.\ 4\% of Normal ones). MSTV is instead
bimodal in the Pathologic class. Splitting Pathologic segments by a label-free
rule explains this: segments with prolonged decelerations ($\mathrm{DP}>0$;
65 segments) have high MSTV, a wide histogram and a low histogram mean (medians
2.7, 113 and 96\,bpm), whereas those without ($\mathrm{DP}=0$; 39 segments)
are flat (medians: MSTV 0.3, Width 18, ALTV 69\%). The two groups deviate from
Normal (MSTV 1.2, Width 68, Mean 134\,bpm, ALTV 0\%) in \emph{opposite}
directions; only ASTV is elevated in both (ASTV $\ge 60$ in 83\% and 95\% of
the two groups vs.\ 16\% of Normal). Event features are zero-inflated, so we
compare presence rates: a prolonged deceleration occurs in 62.5\% of
Pathologic but 4.7\% of Normal segments, and an acceleration in 19.2\% vs.\
67.2\%.

\begin{figure*}[t]
  \centering
  \includegraphics[width=\textwidth]{fig1_variability_contrast.pdf}
  \caption{FHR-variability features in the training portion. Top:
  histograms normalised within each class, so that the 972:104 imbalance does
  not distort the comparison. Bottom: empirical CDFs; the KS statistic is the
  largest vertical gap. ASTV and ALTV are percentages of time; UCI states no
  unit for MSTV and MLTV. Pathologic MSTV is bimodal and 69\% of Pathologic
  ALTV values are zero.}
  \label{fig:variability}
\end{figure*}

\paragraph{Feature separability.}
Table~\ref{tab:sep} reports the two-sample Kolmogorov--Smirnov statistic, a
rank-based AUC, and $k$NN mutual information (MI) with the label. Wasserstein
distances were computed after scaling by the Normal standard deviation, since
raw values are not comparable across units (ASTV: 1.40 SD; Mean: 1.84 SD).
ASTV is the most separable and the most stable feature (KS 0.71, 95\% CI
0.56--0.84). DP and the histogram mean rank next, but their intervals are wide
and their KS drops to 0.41 and 0.37 when the two largest Pathologic
examinations are removed, so they partly reflect a few examinations. A naive
segment-level bootstrap would have given 0.51--0.68 for DP, which understates
this uncertainty. MSTV shows why direction-specific scores mislead for
anomaly detection: its AUC is 0.59, but the AUC of its absolute deviation from
the Normal mean is 0.80. The baseline LB carries no usable signal on its own
(0.37).

\begin{table}[t]
\centering
\small
\setlength{\tabcolsep}{2.8pt}
\begin{tabular}{lrrlccc}
\toprule
 & \multicolumn{2}{c}{Median} & & & & \\
\cmidrule(lr){2-3}
Feature & N & P & KS [95\% CI] & AUC & $|z|$ & MI \\
\midrule
ASTV (\%)     & 40    & 65    & .71 [.56, .84] & .86 & .83 & .13 \\
DP$^{\dagger}$ & 0 & 1.7   & .60 [.29, .78] & .80 & .80 & .13 \\
Mean (bpm)    & 134   & 103.5 & .58 [.29, .77] & .19 & .80 & .14 \\
MLTV          & 8.1   & 0     & .54 [.44, .68] & .17 & .73 & .09 \\
AC$^{\dagger}$ & 2.5 & 0   & .50 [.34, .67] & .23 & .57 & .05 \\
MSTV          & 1.2   & 2.1   & .36 [.31, .59] & .59 & .80 & .13 \\
ALTV (\%)     & 0     & 0     & .22 [.14, .50] & .52 & .63 & .06 \\
LB (bpm)      & 132   & 133   & .22 [.15, .39] & .49 & .37 & .08 \\
\bottomrule
\end{tabular}
\caption{Separability of selected features (training portion; N = Normal,
P = Pathologic). AUC = $P(x_P > x_N)$ with ties counted one half, so values
far from .5 in either direction indicate separation; $|z|$ = AUC of the
absolute $z$-score with respect to the Normal mean and SD; MI in nats (label
entropy 0.32); $\dagger$ events per 1000\,s. CIs resample examinations. All 21 features are in the
supplementary table.}
\label{tab:sep}
\end{table}

\paragraph{Redundancy.}
The histogram family is strongly redundant. $\mathrm{Width}=\mathrm{Max}-\mathrm{Min}$
holds exactly for every training segment, so the covariance of the 21 features
has rank 20. Mode, Mean and Median are nearly interchangeable (Spearman
$\rho=0.90$--$0.97$), and Width, Min and Variance are strongly related
($\rho=-0.89$ and $0.82$). In contrast, the variability features are not
redundant with one another (ASTV--MSTV $\rho=-0.36$; ALTV--MLTV $\rho=-0.02$).
A label-free rule fitted on Normal training segments (complete-linkage
clusters with $|\rho|\ge 0.8$, keep the medoid, drop Width) removes Width,
Min, Mode and Mean, leaving 17 features.

\paragraph{Exploratory hypotheses.}
These observations shape the feature design tested in \S7--\S8. Pathologic
segments form two phenotypes that are visible in different feature families,
which motivates family-level rather than single-feature ablations.
\textbf{H1}: the variability family (ASTV, MSTV, ALTV, MLTV) is necessary;
removing it lowers detection performance, mainly on Pathologic segments
without prolonged decelerations.
\textbf{H2}: deceleration type is informative; replacing DL, DS and DP by
their sum does not help, because light decelerations are common in Normal
segments (45\%) and the sum is less separable than DP alone (KS 0.40 vs.\
0.60). Our initial expectation that aggregated deceleration features would
improve on the raw ones was therefore not supported.
\textbf{H3}: the 17-feature set obtained by the label-free pruning rule
performs on par with all 21 features. Each hypothesis changes only the feature
set, with the data split, model configuration and seeds held fixed, and is
judged on validation PR-AUC.
 inside the ACL template.
% Requires: graphicx, booktabs. Figures live in EDA/figures/ (PDF and PNG).
% Every number below is produced by `python EDA/p2_eda.py` from the
% TRAINING portion only (EDA/outputs/*.csv). Citation keys to add to the
% .bib: huang2021interpretability, ayresdecampos2000sisporto.
% Section references (\S7--\S8) follow the course outline; adjust if needed.
% =====================================================================
\section{Exploratory Feature Analysis}
\label{sec:eda}

All analyses in this section use only the training portion (972 Normal and
104 Pathologic segments from 175 examinations); validation and test data were
not inspected. Labels are used here for description only; the one-class
detectors are fitted on Normal training segments. Because the 104 Pathologic
segments come from only 28 examinations, and two examinations contribute 36 of
them, uncertainty is reported with a cluster bootstrap that resamples whole
examinations (2{,}000 resamples).

\paragraph{Distribution contrast.}
Figure~\ref{fig:variability} compares the four FHR-variability features.
ASTV, the percentage of time with abnormal short-term variability, shifts
upward (median 40 vs.\ 65), and MLTV collapses toward zero (51\% of Pathologic
segments have $\mathrm{MLTV}=0$ vs.\ 4\% of Normal ones). MSTV is instead
bimodal in the Pathologic class. Splitting Pathologic segments by a label-free
rule explains this: segments with prolonged decelerations ($\mathrm{DP}>0$;
65 segments) have high MSTV, a wide histogram and a low histogram mean (medians
2.7, 113 and 96\,bpm), whereas those without ($\mathrm{DP}=0$; 39 segments)
are flat (medians: MSTV 0.3, Width 18, ALTV 69\%). The two groups deviate from
Normal (MSTV 1.2, Width 68, Mean 134\,bpm, ALTV 0\%) in \emph{opposite}
directions; only ASTV is elevated in both (ASTV $\ge 60$ in 83\% and 95\% of
the two groups vs.\ 16\% of Normal). Event features are zero-inflated, so we
compare presence rates: a prolonged deceleration occurs in 62.5\% of
Pathologic but 4.7\% of Normal segments, and an acceleration in 19.2\% vs.\
67.2\%.

\begin{figure*}[t]
  \centering
  \includegraphics[width=\textwidth]{fig1_variability_contrast.pdf}
  \caption{FHR-variability features in the training portion. Top:
  histograms normalised within each class, so that the 972:104 imbalance does
  not distort the comparison. Bottom: empirical CDFs; the KS statistic is the
  largest vertical gap. ASTV and ALTV are percentages of time; UCI states no
  unit for MSTV and MLTV. Pathologic MSTV is bimodal and 69\% of Pathologic
  ALTV values are zero.}
  \label{fig:variability}
\end{figure*}

\paragraph{Feature separability.}
Table~\ref{tab:sep} reports the two-sample Kolmogorov--Smirnov statistic, a
rank-based AUC, and $k$NN mutual information (MI) with the label. Wasserstein
distances were computed after scaling by the Normal standard deviation, since
raw values are not comparable across units (ASTV: 1.40 SD; Mean: 1.84 SD).
ASTV is the most separable and the most stable feature (KS 0.71, 95\% CI
0.56--0.84). DP and the histogram mean rank next, but their intervals are wide
and their KS drops to 0.41 and 0.37 when the two largest Pathologic
examinations are removed, so they partly reflect a few examinations. A naive
segment-level bootstrap would have given 0.51--0.68 for DP, which understates
this uncertainty. MSTV shows why direction-specific scores mislead for
anomaly detection: its AUC is 0.59, but the AUC of its absolute deviation from
the Normal mean is 0.80. The baseline LB carries no usable signal on its own
(0.37).

\begin{table}[t]
\centering
\small
\setlength{\tabcolsep}{2.8pt}
\begin{tabular}{lrrlccc}
\toprule
 & \multicolumn{2}{c}{Median} & & & & \\
\cmidrule(lr){2-3}
Feature & N & P & KS [95\% CI] & AUC & $|z|$ & MI \\
\midrule
ASTV (\%)     & 40    & 65    & .71 [.56, .84] & .86 & .83 & .13 \\
DP$^{\dagger}$ & 0 & 1.7   & .60 [.29, .78] & .80 & .80 & .13 \\
Mean (bpm)    & 134   & 103.5 & .58 [.29, .77] & .19 & .80 & .14 \\
MLTV          & 8.1   & 0     & .54 [.44, .68] & .17 & .73 & .09 \\
AC$^{\dagger}$ & 2.5 & 0   & .50 [.34, .67] & .23 & .57 & .05 \\
MSTV          & 1.2   & 2.1   & .36 [.31, .59] & .59 & .80 & .13 \\
ALTV (\%)     & 0     & 0     & .22 [.14, .50] & .52 & .63 & .06 \\
LB (bpm)      & 132   & 133   & .22 [.15, .39] & .49 & .37 & .08 \\
\bottomrule
\end{tabular}
\caption{Separability of selected features (training portion; N = Normal,
P = Pathologic). AUC = $P(x_P > x_N)$ with ties counted one half, so values
far from .5 in either direction indicate separation; $|z|$ = AUC of the
absolute $z$-score with respect to the Normal mean and SD; MI in nats (label
entropy 0.32); $\dagger$ events per 1000\,s. CIs resample examinations. All 21 features are in the
supplementary table.}
\label{tab:sep}
\end{table}

\paragraph{Redundancy.}
The histogram family is strongly redundant. $\mathrm{Width}=\mathrm{Max}-\mathrm{Min}$
holds exactly for every training segment, so the covariance of the 21 features
has rank 20. Mode, Mean and Median are nearly interchangeable (Spearman
$\rho=0.90$--$0.97$), and Width, Min and Variance are strongly related
($\rho=-0.89$ and $0.82$). In contrast, the variability features are not
redundant with one another (ASTV--MSTV $\rho=-0.36$; ALTV--MLTV $\rho=-0.02$).
A label-free rule fitted on Normal training segments (complete-linkage
clusters with $|\rho|\ge 0.8$, keep the medoid, drop Width) removes Width,
Min, Mode and Mean, leaving 17 features.

\paragraph{Exploratory hypotheses.}
These observations shape the feature design tested in \S7--\S8. Pathologic
segments form two phenotypes that are visible in different feature families,
which motivates family-level rather than single-feature ablations.
\textbf{H1}: the variability family (ASTV, MSTV, ALTV, MLTV) is necessary;
removing it lowers detection performance, mainly on Pathologic segments
without prolonged decelerations.
\textbf{H2}: deceleration type is informative; replacing DL, DS and DP by
their sum does not help, because light decelerations are common in Normal
segments (45\%) and the sum is less separable than DP alone (KS 0.40 vs.\
0.60). Our initial expectation that aggregated deceleration features would
improve on the raw ones was therefore not supported.
\textbf{H3}: the 17-feature set obtained by the label-free pruning rule
performs on par with all 21 features. Each hypothesis changes only the feature
set, with the data split, model configuration and seeds held fixed, and is
judged on validation PR-AUC.
